\documentclass[10pt,journal,compsoc]{IEEEtran}

\usepackage[T1]{fontenc}
\usepackage{microtype}
\usepackage{booktabs}
\usepackage{siunitx}
\usepackage{pgfplots}
\pgfplotsset{compat=1.18}
\usepackage{graphicx}
\usepackage{amsmath}
\usepackage{xcolor}
\usepackage[hidelinks]{hyperref}

\sisetup{detect-weight=true, detect-family=true}

\begin{document}

\title{NumaRing: Topology-Aware Routing for NUMA-Local MPMC Queues,\\and What Broke When We
Optimized It}

\author{Parth~Kumar~Sinha
\IEEEcompsocitemizethanks{
\IEEEcompsocthanksitem P.~K.~Sinha is an independent researcher based in Oxford, United Kingdom.\protect\\
E-mail: sinhaparth555@gmail.com\protect\\
ORCID: \href{https://orcid.org/0009-0002-3120-9301}{0009-0002-3120-9301}
}}

\markboth{}{}

\IEEEtitleabstractindextext{%
\begin{abstract}
Cache-coherence traffic across NUMA sockets, not the choice of lock-free algorithm, is the
dominant cost in a shared-counter multi-producer multi-consumer (MPMC) queue once threads span
more than one socket. We built NumaRing, a queue that gives each NUMA node its own bounded ring
buffer, routes producers and consumers to the ring on their own node, and moves items between
nodes in batches instead of one at a time. We evaluated it against
\texttt{boost::lockfree::queue}, \texttt{moodycamel::ConcurrentQueue}, and a mutex-protected
queue on a real two-socket machine, using a same-host controlled methodology after an initial
cross-host comparison produced misleading numbers that we caught before drawing conclusions from
them. Profiling found that our own topology-routing call cost \SI{183}{\nano\second} per
operation, roughly eleven to twenty times the cost of the ring-buffer operation it gated. Caching
the CPU-to-node lookup cut that to \SI{2.97}{\nano\second} and raised low-contention throughput by
6 to 7$\times$. A second round of fixes to the cross-node work-stealing path removed an
unnecessary shared contention point and added a steal-threshold check, cutting median latency at
32 threads from \SI{1.69}{\milli\second} to \SI{8.4}{\micro\second} (a 202-fold drop) and p99 by
3.1$\times$. We also tried CPU-pause backoff in the ring buffer's compare-and-swap retry loop and
measured it costing 17 to 30 percent of raw throughput under true sustained contention; we
reverted it and report why. Under a 32-vCPU cloud quota that limited us to two NUMA nodes,
throughput at 32 threads stayed well short of the workload's original 120 million operations per
second target. We report the improvements and the ceiling both, with the numbers behind each.
\end{abstract}

\begin{IEEEkeywords}
NUMA, lock-free queue, MPMC, cache coherence, work stealing, tail latency, cache-conscious data
structures
\end{IEEEkeywords}}

\maketitle

\IEEEdisplaynontitleabstractindextext
\IEEEpeerreviewmaketitle

\section{Introduction}

A bounded MPMC queue built on a single shared head and tail counter works well inside one CPU
socket. Every enqueue and dequeue is one compare-and-swap on that counter, and on a single-socket
machine the counter's cache line stays close to every core touching it. That assumption breaks on
a modern server CPU with two or more NUMA sockets. Once producer and consumer threads run on
different sockets, every compare-and-swap on the shared counter forces the cache line holding it
to migrate across the inter-socket interconnect. David et al.\ measured this directly: crossing a
socket boundary costs two to 7.5 times what an intra-socket cache-line transfer costs, and that
gap is a property of the hardware, not of any particular queue implementation
\cite{david2013everything}.

NumaRing addresses this by giving each NUMA node its own queue instead of sharing one. A producer
or consumer routes to the ring buffer on the node it happens to be running on, so the common case
never touches another node's cache lines. When a node's local queue is temporarily full or empty,
the queue moves a batch of items across nodes with one atomic claim instead of claiming items one
at a time, which cuts the number of cross-socket compare-and-swap operations by roughly the batch
size.

Building and evaluating that idea rigorously is most of this paper. Section~\ref{sec:design}
describes the design. Section~\ref{sec:correctness} states the queue's semantics. Section~\ref{sec:eval}
is the evaluation, and it is where the paper earns its keep: a routing-overhead bug we found by
profiling and fixed for a 6 to 7$\times$ low-contention throughput gain, a work-stealing
contention bug we found the same way and fixed for two to three orders of magnitude of tail-latency
improvement at 32 threads, a backoff optimization we tried, measured, and reverted because it cost
real throughput even though it looked like a win in one benchmark, and a methodology mistake
(comparing two different cloud VM instances as if they were the same machine) that we caught
before it became a false claim. Section~\ref{sec:discussion} states plainly what the evaluation
does not show: throughput at 32 threads is nowhere near the project's original target, and nothing
in the data says another round of software fixes closes that gap. That ceiling is real, and a
paper claiming otherwise without the hardware to back it up would not survive review at a systems
venue.

Our contributions are: (1) a hierarchical, NUMA-aware MPMC queue design with batched cross-node
work stealing, built on the bounded MPMC algorithm of Vyukov \cite{vyukov_bounded_mpmc}, itself a
descendant of the Michael--Scott queue \cite{michael1996simple}; (2) a profiling-driven fix to a
per-operation routing cost that we measured, not assumed, at \SI{183}{\nano\second}; (3) a second
profiling-driven fix to work-stealing contention, evaluated with a same-host controlled
methodology after we caught an earlier cross-host measurement mistake; (4) a negative result on
CPU-pause backoff under sustained contention, included because it corrects a common assumption
about backoff always helping; and (5) an honest accounting of what a 32-vCPU cloud quota does and
does not let us claim.

\section{Background and related work}

\textbf{Lock-free MPMC queues.} Michael and Scott's non-blocking queue \cite{michael1996simple}
established the two-pointer, compare-and-swap-based design most later concurrent queues build on.
Vyukov's bounded MPMC queue \cite{vyukov_bounded_mpmc} adapts that idea to a fixed-size ring
buffer: every slot carries its own sequence number, so a producer or consumer only contends on the
one slot it is about to touch, and the shared state is reduced to two atomic position counters. We
use Vyukov's algorithm as the node-local ring buffer inside NumaRing (Section~\ref{sec:ring}).
\texttt{boost::lockfree::queue} \cite{boost_lockfree} and moodycamel's
\texttt{ConcurrentQueue} \cite{desrochers2014moodycamel} are both mature, widely used C++ MPMC
queues; neither is NUMA-topology-aware, and we use both as baselines in
Section~\ref{sec:eval}.

\textbf{NUMA-aware synchronization.} David et al.\ \cite{david2013everything} is the closest
thing to a definitive study of why synchronization stops scaling across sockets: they show that
above a small number of threads, the bottleneck moves from the algorithm to the hardware
cache-coherence protocol, and that this is true across lock-free, lock-based, and hybrid designs
alike. Their result is the premise this paper starts from. We do not claim NumaRing escapes that
hardware ceiling; Section~\ref{sec:discussion} is explicit that it does not, at the scale we could
test.

\textbf{Detecting cross-socket contention.} \texttt{perf c2c} \cite{mario2016c2c} instruments the
Linux performance-counter subsystem to report which cache lines suffer HITM (a load that hits a
line held modified in another core's cache), the direct hardware signature of the problem this
paper is about. We planned to use it. Section~\ref{sec:eval} explains why the cloud environment we
tested on made that impossible, and what evidence we used instead.

\section{Design}
\label{sec:design}

NumaRing is a header-only C++20 library. Its top-level type, \texttt{Queue<T>}, owns one
\texttt{NodeLocalRingBuffer<T>} per NUMA node reported by the host (Section~\ref{sec:topology}) and
adds a batched cross-node work-stealing layer on top (Section~\ref{sec:worksteal}). On a
single-node or non-NUMA host it degrades to exactly one ring buffer, so the design costs nothing on
hardware that does not need it.

\subsection{Node-local ring buffer}
\label{sec:ring}

Each node's queue is a bounded, lock-free ring buffer implementing Vyukov's algorithm
\cite{vyukov_bounded_mpmc}. Every slot holds a value and its own atomic sequence number. An
enqueue reads the current tail position, checks the target slot's sequence number to confirm the
slot is free, and claims the slot with one compare-and-swap on the tail counter; a dequeue is the
mirror image on the head counter. Because the sequence number lives in the slot itself rather than
in the shared counter, two threads only actually contend when they land on the same slot in the
same generation, which under normal load is rare. This is the property that makes the algorithm
fast on a single socket; it says nothing about what happens across sockets, which is the problem
NumaRing's outer layer exists to address.

The ring's backing storage is allocated with \texttt{numa\_alloc\_onnode()} on the node the queue
is constructed for, so the memory a node's threads touch on the fast path is local DRAM, not a
cross-socket fetch. Both position counters and the per-slot sequence numbers are padded to
\SI{128}{\byte} (two cache lines) to prevent the L2 spatial prefetcher from pulling adjacent,
independently-written lines into the same coherence unit, which would otherwise reintroduce false
sharing that per-slot sequencing was supposed to eliminate.

\subsection{Topology-aware routing}
\label{sec:topology}

A producer or consumer thread finds its own ring buffer by asking which NUMA node it is currently
running on (\texttt{sched\_getcpu()} plus a CPU-to-node lookup) and indexing into the per-node
array of ring buffers with the result. This call happens on every single \texttt{try\_enqueue} and
\texttt{try\_dequeue}, so its cost is not incidental; Section~\ref{sec:eval-routing} measures it
directly and reports what we did about it.

\subsection{Batched cross-node work stealing}
\label{sec:worksteal}

When a thread's local ring buffer is full (on enqueue) or empty (on dequeue), \texttt{Queue<T>}
does not fail immediately. It first tries to move a batch of items between the local ring and
another node's ring, then retries the local operation. A batch transfer claims a contiguous run of
slots with one compare-and-swap on each side instead of one compare-and-swap per item, so moving
16 items this way costs 2 cross-node atomic operations instead of 32, a 93.75\% reduction on the
line every other core sharing that ring buffer is also contending for. If the batch transfer
cannot move anything (every other node's ring is also full or also empty), \texttt{Queue<T>} falls
back to a single-item attempt against every other node in turn before reporting failure to the
caller, so correctness never depends on the fast batched path succeeding.

Section~\ref{sec:eval-worksteal} describes two problems we found in this layer once we started
measuring it at higher thread counts, and the fixes for both.

\section{Correctness}
\label{sec:correctness}

Within one node's ring buffer, NumaRing inherits Vyukov's linearizability and FIFO-per-slot
guarantees directly: the per-slot sequence number is the single point of agreement between
producers and consumers, and every enqueue or dequeue that returns \texttt{true} corresponds to
exactly one atomic transition of that number. A ring buffer's \texttt{try\_enqueue} either
succeeds and stores the value, or fails and leaves the caller's value completely untouched, which
\texttt{Queue<T>}'s overflow handling depends on: a failed local enqueue can be retried against
another node without a defensive copy.

Across nodes, \texttt{Queue<T>} does not offer a single global FIFO order, and does not claim to.
An item enqueued on node~0 and later stolen to node~1 can be dequeued before an item enqueued
slightly earlier that stayed on node~0 the whole time. This is a deliberate trade: a queue that
enforced strict global ordering across sockets would need the same shared cross-socket counter this
design exists to avoid. What NumaRing does guarantee is that a batch transfer is atomic from every
other thread's point of view (a concurrent observer sees either the whole batch moved or none of
it, never a partial batch), and that the fallback single-item path preserves the same per-slot
sequencing guarantee the local ring buffer provides.

\section{Evaluation}
\label{sec:eval}

\subsection{Setup and methodology}
\label{sec:eval-setup}

All numbers in this section come from a throwaway Google Compute Engine \texttt{n2-standard-32}
instance running Ubuntu 24.04, built fresh for each measurement round and deleted afterward. The
instance exposes 32 vCPUs split across two NUMA nodes, 16 vCPUs each, with a measured node distance
of 20 (versus 10 within a node), confirmed with \texttt{numactl -{}-hardware} before every run. This
is a hard ceiling, not a choice: the GCP project used for this work has a 32-vCPU total quota
across every region, and every machine type we found that reliably exposes more than two NUMA
nodes needs more vCPUs than that. We built with GCC~13 at \texttt{-O3 -march=native}, and compared
against Boost's \texttt{lockfree::queue} (built with \texttt{fixed\_sized<true>} for a true bounded
ring, matching NumaRing's own fixed capacity) and moodycamel's \texttt{ConcurrentQueue}, both at
their latest stable release, plus a queue built from \texttt{std::mutex} and \texttt{std::deque} as
the traditional-locking baseline.

The workload pins one producer and one consumer to each NUMA node and has every producer enqueue a
monotonic timestamp as its payload; a consumer's latency sample is the wall-clock time between that
timestamp being written and the same value being dequeued by whichever thread gets it, which
measures real cross-thread queueing delay rather than a same-thread round trip. All four
implementations run the identical workload through the identical harness, with thread affinity set
explicitly (\texttt{numa\_run\_on\_node()}) rather than left to the scheduler, because a locality
result is not meaningful if the thread-to-node mapping the queue is routing against is not itself
stable.

Our first attempt at measuring the routing and work-stealing fixes compared a pre-fix build on one
GCE instance against a post-fix build on a second, freshly created instance. Every implementation's
numbers moved on the second instance, including \texttt{boost::lockfree::queue} and
\texttt{moodycamel::ConcurrentQueue}, neither of which our changes touch. A code change confined to
\texttt{numaring::Queue} cannot make an unrelated library slower. The only explanation is that the
two GCE instances landed on different physical hosts with different noisy-neighbor conditions, a
known pitfall in cloud benchmarking. We caught this by checking whether libraries we had not
touched moved too, before writing any comparison number down, and fixed it by checking out the
pre-fix commit into a second build directory on the \emph{same} instance (via a \texttt{git
worktree}) and running both builds back to back. Every result below that compares a before and an
after state used that same-host method.

\subsection{Routing overhead}
\label{sec:eval-routing}

Profiling \texttt{current\_node()}, the call every \texttt{try\_enqueue} and \texttt{try\_dequeue}
makes to find its caller's local ring buffer, isolated the cost with two targeted microbenchmarks:
\texttt{sched\_getcpu()} alone costs \SI{2.3}{}--\SI{3.0}{\nano\second} (it resolves through a vDSO
call, not a real syscall trap), while \texttt{numa\_node\_of\_cpu()} alone costs
\SI{130}{}--\SI{131}{\nano\second}, almost the entire measured cost of \texttt{current\_node()}. It
is a bitmask scan over every NUMA node that libnuma redoes from scratch on every call, despite the
answer for a given CPU never changing for the life of the process. We replaced it with a
lazily-built, process-wide lookup table from CPU to node, computed once and read thereafter;
\texttt{sched\_getcpu()} is still called fresh on every operation, so no staleness was introduced,
only the static part of the lookup is cached.

\begin{table}[t]
\caption{Routing-lookup cost, measured directly on real 2-node NUMA hardware.}
\label{tab:routing}
\centering\small
\begin{tabular}{lrrr}
\toprule
Metric & Before & After & Factor \\
\midrule
\texttt{current\_node()} & \SI{183}{ns} & \SI{2.97}{ns} & 61.6$\times$ \\
\texttt{numa\_node\_of\_cpu()} & \SI{130}{ns} & \SI{0.89}{ns} & 146$\times$ \\
\bottomrule
\end{tabular}
\end{table}

For context, a single-threaded enqueue-dequeue round trip through the underlying ring buffer, with
no routing call at all, costs \SI{16.2}{\nano\second}. Before this fix, routing alone cost more
than eleven times that; after, it costs less than a fifth of it.

\subsection{Work-stealing contention}
\label{sec:eval-worksteal}

The routing fix alone left a puzzle: it delivered a large gain at 2--8 threads and almost none at
16--32, where tail latency stayed in the millisecond range. The shape of that result, a fast median
next to a slow p99 and p99.9, is the signature of a shared contended resource rather than raw
per-operation cost, so we looked for one in the work-stealing path and found two.

First, candidate selection for which other node to steal from or spill to went through a single
\texttt{std::atomic<std::size\_t>} that every thread incremented on every overflow or underflow
event. Under saturation, when most producers on most nodes are overflowing at once, that counter
becomes a cache line contended by every thread on the fast-failure path, for no benefit beyond
spreading out which candidate node gets tried first. We replaced it with per-thread state, seeded
once per thread and advanced independently, which needs no cross-thread synchronization at all.

Second, a steal attempt would proceed against any candidate node with at least one item available,
which under saturation meant several threads simultaneously racing to steal the last few items from
a nearly empty source, most of their compare-and-swap attempts doomed to fail and retry. We added a
cheap, deliberately approximate size check (two relaxed atomic loads, never used for a correctness
decision) and gated a steal attempt on the source having at least twice the batch size available. A
skipped attempt falls back to the existing single-item cross-node path, so this only changes which
code path handles the low-occupancy case, never correctness.

\begin{table}[t]
\caption{\texttt{numaring::Queue} throughput and median latency, same host, before and after both
routing and work-stealing fixes.}
\label{tab:worksteal}
\centering\small
\begin{tabular}{lrrrr}
\toprule
Threads & \multicolumn{2}{c}{Throughput (M\,ops/s)} & \multicolumn{2}{c}{p50 latency} \\
 & Before & After & Before & After \\
\midrule
2  & 1.73  & 12.35 & \SI{3.48}{ms} & \SI{218}{\micro\second} \\
32 & 0.260 & 0.321 & \SI{1.69}{ms} & \SI{8.4}{\micro\second} \\
\bottomrule
\end{tabular}
\end{table}

Figure~\ref{fig:scaling} plots throughput against thread count for the fixed build against all
three baselines, on the same host, same run. Figure~\ref{fig:tail} plots p50/p95/p99/p99.9 latency
for the same four implementations at 32 threads.

\begin{figure}[t]
\centering
\begin{tikzpicture}
\begin{axis}[
  width=\linewidth, height=5.2cm,
  xlabel={Threads}, ylabel={Throughput (M items/s)},
  xmode=log, log basis x=2, ymode=log,
  xtick={2,4,8,16,32}, xticklabels={2,4,8,16,32},
  legend style={font=\scriptsize, at={(0.98,0.98)}, anchor=north east},
  grid=both, grid style={line width=.1pt, draw=gray!25},
]
\addplot[mark=*, thick] table[x=threads,y=throughput] {data/postfix_numaring.dat};
\addlegendentry{NumaRing}
\addplot[mark=square*, thick, color=orange] table[x=threads,y=throughput] {data/postfix_boostlockfree.dat};
\addlegendentry{Boost.Lockfree}
\addplot[mark=triangle*, thick, color=teal] table[x=threads,y=throughput] {data/postfix_moodycamel.dat};
\addlegendentry{moodycamel}
\addplot[mark=diamond*, thick, color=gray] table[x=threads,y=throughput] {data/postfix_mutex.dat};
\addlegendentry{Mutex+deque}
\end{axis}
\end{tikzpicture}
\caption{Throughput vs.\ thread count, same host, current build. All four implementations run the
identical timestamp-payload workload through the identical harness.}
\label{fig:scaling}
\end{figure}

\begin{figure}[t]
\centering
\begin{tikzpicture}
\begin{axis}[
  width=\linewidth, height=5.2cm,
  ybar, bar width=4pt,
  symbolic x coords={p50,p95,p99,p999},
  xtick=data,
  ylabel={Latency ($\mu$s, log scale)}, ymode=log,
  legend style={font=\scriptsize, at={(0.02,0.98)}, anchor=north west},
  grid=major, grid style={line width=.1pt, draw=gray!25},
  enlarge x limits=0.2,
]
\addplot[fill=blue!60] table[x=metric,y=numaring] {data/tail_at_32.dat};
\addplot[fill=orange!70] table[x=metric,y=boost] {data/tail_at_32.dat};
\addplot[fill=teal!60] table[x=metric,y=moodycamel] {data/tail_at_32.dat};
\addplot[fill=gray!50] table[x=metric,y=mutex] {data/tail_at_32.dat};
\legend{NumaRing, Boost.Lockfree, moodycamel, Mutex+deque}
\end{axis}
\end{tikzpicture}
\caption{p50/p95/p99/p99.9 latency at 32 threads, same host, current build. See
Table~\ref{tab:tail32} for the exact values plotted.}
\label{fig:tail}
\end{figure}

\begin{table}[t]
\caption{Latency percentiles at 32 threads, same host, current build.}
\label{tab:tail32}
\centering\small
\begin{tabular}{lrrrr}
\toprule
Implementation & p50 & p95 & p99 & p99.9 \\
\midrule
NumaRing        & \SI{8.4}{\micro\second}   & \SI{799}{\micro\second}  & \SI{1.23}{ms} & \SI{1.69}{ms} \\
Boost.Lockfree  & \SI{3.99}{ms}  & \SI{4.40}{ms}  & \SI{4.45}{ms} & \SI{4.51}{ms} \\
moodycamel      & \SI{4.42}{\micro\second}   & \SI{2.50}{ms}  & \SI{3.17}{ms} & \SI{4.19}{ms} \\
Mutex+deque     & \SI{43.3}{\micro\second}   & \SI{160}{\micro\second}  & \SI{227}{\micro\second} & \SI{290}{\micro\second} \\
\bottomrule
\end{tabular}
\end{table}

Table~\ref{tab:tail32} is worth reading carefully rather than skimming for a winner. NumaRing's
median is the lowest of the four at 32 threads, but its tail is not: p99 and p99.9 land between
Boost's and moodycamel's, both well above the mutex baseline's tail at this specific thread count
and workload. We are not going to round that up. Section~\ref{sec:discussion} says what we think
is actually going on.

\subsection{NUMA memory locality}

\texttt{numastat}, sampled before and after a 15-second sustained run at 32 threads, showed
\texttt{numa\_hit}/\texttt{numa\_miss} ratios of 100\% on both nodes, with zero misses recorded on
either. A per-process snapshot taken mid-run (\texttt{numastat -p}) showed private memory split
across both nodes rather than concentrated on one, consistent with one \texttt{NodeLocalRingBuffer}
correctly landing on each node's own memory via \texttt{numa\_alloc\_onnode()}. Because
\texttt{NodeLocalRingBuffer} allocates its backing storage once at construction and never again
during steady-state operation, the system-wide counters move only slightly during a run; that is
expected given the design, not a measurement gap.

\subsection{A negative result: CPU-pause backoff}
\label{sec:eval-backoff}

Standard advice for a contended compare-and-swap retry loop is to back off: pause briefly before
retrying, so a losing thread does not immediately re-contend the same cache line. We implemented
two variants in the ring buffer's own retry loops, a flat \texttt{\_\_builtin\_ia32\_pause}-per-retry
version and a version that grows the pause geometrically with consecutive failures up to a bound,
and measured both against an uninstrumented, fixed-duration sustained-load driver at 32 threads
that has no per-operation clock call diluting the measurement.

\begin{table}[t]
\caption{\texttt{numaring\_sustained\_load} throughput (M ops/s) at 32 threads, same host,
15-second runs, three repeats per variant.}
\label{tab:backoff}
\centering\small
\begin{tabular}{lrrrr}
\toprule
Variant & Run 1 & Run 2 & Run 3 & Mean \\
\midrule
Pre-fix (no backoff)     & 8.65 & 9.41  & 9.02  & 9.03 \\
Flat pause               & 7.44 & 8.08  & 6.95  & 7.49 \\
Exponential backoff      & 7.25 & 5.44  & 6.47  & 6.39 \\
Work-stealing fix, no backoff & 9.82 & 8.24 & 11.03 & \textbf{9.70} \\
\bottomrule
\end{tabular}
\end{table}

Both backoff variants reduced raw throughput, and the more sophisticated one reduced it further,
not less. Under this driver's true wall-to-wall contention, with no pacing between a thread's own
retries, the fastest path back to success for a losing compare-and-swap attempt is to retry
immediately; the atomic read-modify-write itself already serializes access through the cache
coherence protocol, so spending cycles on a pause instruction is pure loss, and a growing pause
loses more. We had measured the same backoff variants \emph{helping} tail latency in the
instrumented benchmark from Section~\ref{sec:eval-worksteal}, which uses a real clock call around
every operation; our best explanation for the discrepancy is that the clock call itself inserts
wall-clock spacing between a thread's successive retries that the uninstrumented driver does not
have, and that spacing is enough to make a fixed pause look free when it is not. We reverted both
backoff variants and report this result because it corrects an assumption (that backoff helps a
retry loop) that does not hold under every workload shape, and because a paper that only reports
what worked would be reporting half the evidence.

\subsection{What we could not measure}

\texttt{perf c2c} needs the host to expose a hardware performance-monitoring unit with support for
precise memory-load sampling. \texttt{perf stat -e cycles,instructions} on our test instance
reported every counter, including plain cycle counting, as \texttt{<not supported>}; the kernel log
confirmed the reason directly: \texttt{"Performance Events: unsupported CPU family 6 model 85 no
PMU driver, software events only."} GCP has a
\texttt{-{}-performance-monitoring-unit=enhanced}
flag for exactly this, but it was rejected for every machine type we could reach under our 32-vCPU
quota, including both \texttt{n2-standard-32} and \texttt{c3-standard-22}. We are reporting this as
a limitation rather than omitting it, because a paper that is honest about what its evaluation shows
should also be honest about what it does not.

\section{Discussion and limitations}
\label{sec:discussion}

The evaluation supports two conclusions that both matter and point in different directions.

The routing fix and the work-stealing fix are real, and the tail-latency numbers behind them are
not close calls: a 202-fold drop in median latency and a 3.1-fold drop in p99 at 32 threads, on the
same host, measured before and after with nothing else changed. Anyone building a NUMA-local queue
on a hot path that queries thread affinity per operation should expect the same class of bug we
found, because the cost of a topology query is easy to assume is negligible and expensive to
verify is not.

The throughput ceiling is also real, and no amount of software tuning we tried moved it much. At 32
threads across two sockets, NumaRing's raw throughput sits in the low hundreds of thousands of
operations per second in the instrumented benchmark and around ten million in the uninstrumented
driver, both far short of the workload's original 120 million operations per second target. We do
not think a third round of the same kind of fix closes that gap. The two fixes in this paper
targeted software overhead this project's own code introduced (a redundant lookup, an unnecessary
shared atomic); once those are gone, what remains is the same hardware limit the motivation in
Section~\ref{sec:eval-setup} describes: two sockets connected by an interconnect with fixed
bandwidth and fixed latency, and a workload that deliberately routes half its traffic across that
interconnect. Closing that gap, if it can be closed at all with this design, needs either a larger
NUMA topology than two sockets to test against, or a different approach to how the shared position
counters are contended, and we do not have evidence for either from a 32-vCPU cloud quota.

The evaluation's other limits are worth stating plainly rather than leaving implicit. Every number
in Section~\ref{sec:eval} comes from one machine type (\texttt{n2-standard-32}) in one cloud
region, and while we controlled for host-to-host variance within that constraint, we have not
tested a different NUMA topology, a different interconnect, or bare-metal hardware. The workload is
a synthetic, symmetric producer-consumer benchmark with a fixed capacity and batch size; it is not
a trace from a real application, and different workload shapes could move the work-stealing fix's
benefit up or down. \texttt{perf c2c} confirmation of the cross-socket cache-invalidation reduction
this design is meant to deliver remains unmeasured, for the environmental reason given in
Section~\ref{sec:eval}, not because we chose not to look.

\section{Conclusion}

We built a NUMA-aware MPMC queue, found and fixed two real performance bugs by profiling rather
than guessing, tried an optimization that looked correct and measured it costing throughput instead
of saving it, caught a benchmarking mistake before it became a false claim, and are reporting a
throughput target we did not hit alongside the latency improvements we did. We think that is a more
useful contribution than a set of numbers tuned to look better than the hardware we tested on
actually supports.

\bibliographystyle{IEEEtran}
\bibliography{references}

\end{document}
